\documentclass[11pt]{article}
\usepackage[margin=1in, headheight=14pt]{geometry}
\usepackage{amsmath, amssymb}
\usepackage{hyperref}
\usepackage{enumitem}
\usepackage{algorithm}
\usepackage{algpseudocode}
\usepackage{fancyhdr}
\usepackage{tcolorbox}
\usepackage{microtype}
\usepackage{tikz}
\usepackage{pgfplots}
\pgfplotsset{compat=1.18}
\usetikzlibrary{arrows.meta, positioning, calc, shapes.geometric}

\input{notation.tex}
\input{zk-notation.tex}
\input{environments.tex}
\input{lecture-macros.tex}

\pagestyle{fancy}
\fancyhf{}
\lhead{EECS 498/CSE 598: Zero-Knowledge Proofs}
\rhead{Winter 2026}
\lfoot{Lecture 18: Memory Fingerprinting, Lookup Arguments, and Plookup}
\rfoot{\thepage}
\renewcommand{\headrulewidth}{0.4pt}
\renewcommand{\footrulewidth}{0.4pt}

% Lecture-specific macros
\tcbuselibrary{skins, breakable}

\newtcolorbox{highlight}[1][]{
  colback=blue!5!white,
  colframe=blue!60!black,
  fonttitle=\bfseries,
  #1
}

\newtcolorbox{graybox}[1][]{
  colback=gray!8,
  colframe=gray!50,
  fonttitle=\bfseries,
  #1
}

\newcommand{\Plookup}{\ensuremath{\mathsf{Plookup}}}

\begin{document}

%-----------------------------------------------------------------------
% Title Block
%-----------------------------------------------------------------------
\begin{center}
  {\LARGE \textbf{EECS 498 / CSE 598: Zero-Knowledge Proofs}}\\[6pt]
  {\large Winter 2026}\\[4pt]
  {\large \textbf{Lecture 18: Memory Fingerprinting, Lookup Arguments, and Plookup}}\\[4pt]
  {\large Instructor: Paul Grubbs \quad
    \href{mailto:paulgrub@umich.edu}{\texttt{paulgrub@umich.edu}} \quad
    Beyster 4709}
\end{center}

\vspace{0.5em}
\hrule
\vspace{1em}

%-----------------------------------------------------------------------
% Agenda
%-----------------------------------------------------------------------
\section*{Agenda}
\begin{itemize}
  \item Memory fingerprinting
  \item Multiset equality via polynomial identity testing
  \item Randomized checks inside CPU-model frontends
  \item Lookup arguments
  \item Simplified $\mathsf{Plookup}$
  \item Modern CPU-model frontends
\end{itemize}

\hrule
\vspace{1em}

%-----------------------------------------------------------------------
\section{Memory Fingerprinting}
%-----------------------------------------------------------------------

Memory fingerprinting is an alternative to Merkle-tree-based authenticated
memory for proving memory consistency in CPU-model frontends.  The dominant
advantage is not asymptotic avoidance of processing memory accesses, but a
substantial improvement in the per-access constants.

\subsection{Randomized Circuits}

The fingerprinting construction uses a randomized verification step.  The
relevant abstraction is a randomized circuit
\[
  C(x,w;r),
\]
where $x$ is the public input, $w$ is the witness, and $r$ is verifier-chosen
randomness.  Soundness is taken over the random choice of $r$.

\begin{definition}[Randomized circuit]
A randomized circuit is a circuit with public input $x$, witness $w$, and
verifier challenge $r$, together with a soundness guarantee over the choice of
$r$.
\end{definition}

A randomized circuit is the algebraic analogue of a public-coin
Merlin--Arthur verification procedure.

\begin{highlight}[title={Randomized Memory-Fingerprinting Pattern}]
For a fixed public input $x$:
\begin{enumerate}
  \item Commit to the portion of the witness that does not depend on the
        challenge and obtain a pre-challenge commitment $c_p$.
  \item Sample a verifier challenge $r \getsr \F$.
  \item Compute the remaining witness data and a proof $\pi$ that
        $C(x,w;r)=1$.
  \item Verify $\pi$ relative to $(x,c_p,r)$.
\end{enumerate}
\end{highlight}

When the protocol is public coin, this interaction can be compiled away with
Fiat--Shamir by seeding the transcript with the commitment $c_p$.

\subsection{Restricted Memory Transcripts}

The memory-fingerprinting argument works with memory-access tuples of the form
\[
  (\tau,\ a,\ v,\ \tau_{\mathrm{prev}}),
\]
where:
\begin{itemize}
  \item $\tau$ is the current timestamp,
  \item $a$ is the accessed location,
  \item $v$ is the value returned or written,
  \item $\tau_{\mathrm{prev}}$ is version information identifying the previous write
        associated with that value.
\end{itemize}

The reduction to multiset equality uses the following structural restrictions
on the transcript.

\begin{enumerate}
  \item Add an initial write to every memory location.
  \item Before every write, insert a dummy read of the same location.
  \item After every read, insert a write-back of the value that was just read.
  \item On every write, store the current timestamp together with the value.
  \item After program termination, read every memory location once.
  \item Reject if any read returns a timestamp larger than the current
        timestamp.
\end{enumerate}

If memory has $s$ locations, the initial writes can be indexed by timestamps
$0,\ldots,s-1$.

The timestamp component is a version tag.  Without it, an adversary could try
to justify a later read using an older value that had once been stored at the
same address.  The reject condition prevents ``reads from the future'' and is
part of the versioning invariant.

\subsection{Reduction to Multiset Equality}

\begin{proposition}[Memory consistency from multiset equality]
Consider a memory transcript satisfying the restrictions above, and assume that
the reject condition is never triggered.  Then the transcript is consistent,
meaning that every read returns the most recently written value, if and only if
the sequence of read operations and the sequence of write operations are equal
as multisets.
\end{proposition}

\begin{proof}[Proof sketch]
Each read contributes a tuple whose value and version information must match a
previous write.  If the transcript is consistent, every read tuple is matched
by exactly one write tuple, and the two multisets coincide.

Conversely, if the multisets coincide and the versioning invariant is enforced,
then every read must correspond to some previous write of the same versioned
value.  The version tag rules out reordering attacks in which a stale value is
reused after a later write to the same address.
\end{proof}

\subsection{Encoding Tuples as Field Elements}

The multiset-equality protocol works over field elements.  The simplest case is
the injective one: if the tuple
\[
  (\tau,a,v,\tau_{\mathrm{prev}})
\]
fits into a single field element by concatenation, then tuple equality is
exactly equality of the encoded field elements.

If the tuple is too large to concatenate directly, a randomized compression can
be used instead.  The verifier supplies random coefficients and the tuple is
compressed into a single field element with a collision probability bounded in
terms of the field size.  The injective concatenation case is used below.

%-----------------------------------------------------------------------
\section{Multiset Equality via Polynomial Identity Testing}
%-----------------------------------------------------------------------

Let
\[
  A = (a_1,\ldots,a_t)
  \qquad\text{and}\qquad
  B = (b_1,\ldots,b_t)
\]
be multisets of elements of $\F$.

\begin{definition}[Characteristic polynomial of a multiset]
Define
\[
  P_A(X) = \prod_{i=1}^{t}(a_i - X)
  \qquad\text{and}\qquad
  P_B(X) = \prod_{i=1}^{t}(b_i - X).
\]
\end{definition}

The multisets are equal exactly when the two polynomials are equal.

\begin{protocol}[Randomized multiset-equality check]
Given oracle access to $P_A$ and $P_B$:
\begin{enumerate}
  \item the verifier samples $r \getsr \F$;
  \item the prover returns $P_A(r)$ and $P_B(r)$;
  \item the verifier accepts iff $P_A(r)=P_B(r)$.
\end{enumerate}
\end{protocol}

\begin{proposition}[Soundness]
If $A \neq B$ as multisets, then the multiset-equality test above accepts with
probability at most
\[
  \frac{t}{|\F|}.
\]
\end{proposition}

\begin{proof}
If $A \neq B$, then $P_A - P_B$ is a nonzero polynomial of degree at most $t$.
A nonzero polynomial of degree at most $t$ has at most $t$ roots in $\F$.
Hence
\[
  \Pr_{r \getsr \F}[P_A(r)=P_B(r)]
  =
  \Pr_{r \getsr \F}[(P_A-P_B)(r)=0]
  \le \frac{t}{|\F|}.
\]
\end{proof}

Under injective tuple encoding, this polynomial test can be applied directly to
the multisets of memory-access tuples.

%-----------------------------------------------------------------------
\section{Randomized Checks Inside CPU-Model Frontends}
%-----------------------------------------------------------------------

The CPU-step relation itself remains deterministic.  The randomized part enters
only in the memory-consistency check.  A direct frontend architecture first
computes the deterministic execution trace, uses local step checks on adjacent
states, and extracts the memory transcript that will later be checked
separately.

The witness then splits naturally into two portions:
\[
  Z_1 = (x,\ s_1,\ldots,s_T,\ \text{deterministic step and preprocessing data}),
\]
\[
  Z_2 = (\text{post-challenge fingerprint evaluations and auxiliary data}).
\]

The state trace can be computed before $r$ is known, because the execution of
the program itself is deterministic once the input is fixed.  The fingerprint
subprotocol is the only part that needs the verifier's challenge.

The pre-challenge witness consists of the full execution trace together with
the deterministic preprocessing on the memory transcript.  Only the final
fingerprint evaluations depend on the verifier challenge.  The local step
checks and the initial memory checks therefore remain in $Z_1$, while the
challenge-dependent memory fingerprint forms the relevant portion of $Z_2$.

In an $\algo{R1CS}$ implementation, the deterministic step checks appear as the
usual local transition checks on adjacent states.  The randomized memory check
is attached only after the verifier challenge is fixed.  Since the memory
fingerprinting protocol is public coin, the extra interaction can again be
removed by Fiat--Shamir, seeded by the commitment to $Z_1$.

%-----------------------------------------------------------------------
\section{Lookup Arguments}
%-----------------------------------------------------------------------

Most of the complexity of the full memory-fingerprinting problem comes from
writes.  In a read-only setting, timestamps and versioning disappear, and the
problem becomes much simpler.

Fix a public immutable table $T = (T[0], T[1], \ldots)$.  The witness now
contains a sequence of claimed index-value pairs
\[
  (a_1,v_1),\ldots,(a_t,v_t),
\]
and the task is to show that $T[a_i] = v_i$ for every $i$.

After concatenating or randomly compressing each pair into a single field
element, the problem becomes the following: given a public table and a witness
sequence
\[
  f_1,\ldots,f_N,
\]
prove that every $f_i$ is contained in the public table.

This is the setting for lookup arguments.

%-----------------------------------------------------------------------
\section{Simplified \texorpdfstring{\Plookup}{Plookup}}
%-----------------------------------------------------------------------

Write the public lookup table as $S = (S_1,\ldots,S_B)$, the claimed lookup
values as $F = (F_1,\ldots,F_N)$, and the witness as
$W = (W_1,\ldots,W_N)$, where $W$ is claimed to be the sorted version of $F$.

\begin{graybox}[title={Simplified \Plookup\ Local Conditions}]
Public data:
\[
  S = (S_1,\ldots,S_B), \qquad F = (F_1,\ldots,F_N).
\]
Witness:
\[
  W = (W_1,\ldots,W_N),
\]
where $W$ is claimed to be the sorted version of $F$.

Assumptions:
\begin{enumerate}
  \item every table element $S_i$ is looked up at least once;
  \item $S_1,\ldots,S_B$ is a contiguous interval of integers.
\end{enumerate}

Checks:
\begin{enumerate}
  \item $W$ is a permutation of $F$;
  \item $W_1 = S_1$;
  \item $W_N = S_B$;
  \item for every $i=2,\ldots,N$,
        \[
          (W_i - W_{i-1})\bigl(W_i - (W_{i-1}+1)\bigr) = 0.
        \]
\end{enumerate}
\end{graybox}

The last constraint says that adjacent entries of $W$ differ by either $0$ or
$1$.  Because the table is a contiguous interval and the endpoints are fixed,
these local constraints imply that the sorted witness never leaves the table
and cannot skip any table value.

\begin{proposition}[Correctness of the simplified check]
Under the two simplifying assumptions above, if $W$ satisfies the four checks,
then every $F_i$ belongs to the public table $S$.
\end{proposition}

\begin{proof}[Proof sketch]
The permutation check implies that $W$ and $F$ contain the same elements with
the same multiplicities.  The endpoint checks and the local constraints force
$W$ to be a nondecreasing sequence whose adjacent differences are always $0$ or
$1$, beginning at $S_1$ and ending at $S_B$.  Every element of $W$ therefore
lies in the interval defined by the table, and hence every element of $F$ does
as well.
\end{proof}

\subsection{Cost}

The permutation check can be implemented with a polynomial-identity or
permutation argument, and therefore costs $O(N)$ constraints.  The endpoint
checks are constant cost, and the local relation
\[
  (W_i - W_{i-1})\bigl(W_i - (W_{i-1}+1)\bigr) = 0
\]
costs only a constant number of constraints per position.  The overall cost is
therefore linear in $N$.

\begin{remark}
The argument reduces a global membership statement to a local sortedness
relation plus a permutation check.
\end{remark}

%-----------------------------------------------------------------------
\section{Modern CPU-Model Frontends}
%-----------------------------------------------------------------------

CPU-model frontends provide an end-to-end path from ordinary compiled programs
to zero-knowledge proofs of execution.  At the level of abstraction, a frontend
for a sufficiently expressive instruction set removes the need to design a
separate ASIC-style proving language for every application.

In principle, a proof system for a real instruction set architecture would make
it possible to compile an ordinary program to assembly and then prove facts
about the resulting execution trace directly.  In practice, architectures such
as x86 carry substantial legacy complexity, including features such as
binary-coded-decimal instructions that are awkward for proof systems.  Current
systems therefore tend to favor cleaner architectures such as RISC-V.

Read-only lookup arguments fit naturally into this broader picture because many
useful CPU tables are immutable, and immutable tables admit substantially
simpler consistency arguments than fully mutable memory.

\begin{remark}
The structural split is stable across modern designs: local transition checks
handle CPU execution, memory fingerprinting handles mutable memory, and lookup
arguments handle immutable tables.
\end{remark}

\bigskip
\hrule
\bigskip

\begin{thebibliography}{9}

\bibitem{GabizonWilliamson20}
  A.~Gabizon and Z.~J. Williamson,
  ``plookup: A simplified polynomial protocol for lookup tables,''
  \emph{IACR ePrint 2020/315}, 2020.

\end{thebibliography}

\end{document}
